% The next line makes it so it compiles the main file if I hit the compile button
% when editing this file in TeXworks.
% !TEX root = ../ActiveIntroToDiscreteMathAndAlgorithms.tex
%%%%%%%%%%%%%%%%%%%%%%%%%%%%%%%%%%%%
\section{Problems}

\begin{pro}
Prove Theorem \ref{thm:theta}.
\end{pro}

\begin{pro}
%This is done in an earlier exercise, but oh well.
$\Theta$ can be thought of as a relation on the set of positive functions, 
where $(f,g)\in\Theta$ iff $f(n)=\Theta(g(n))$.
Prove that $\Theta$ is an equivalence relation. 
\end{pro}

\begin{pro}
Rank the following functions in increasing rate of growth.  
Indicate if two or more functions have the same growth rate.
$$     
x!, \,  
x^3,\, 
x^2\log x,\, \,
x,\, \,     
x^{\log_2 3},\, \,   
\sqrt{x},\, \,   
3^x,\, \,   
x \log x,\, \,  
x^2,\, \,    
x^x,\, \, 
x^{3/2},\, \, 
x^{\log_3 7},\, \,  
x\log(x^2),\, \, 
x\log(\log(x)),\, \, 
\left(\frac{3}{2}\right)^x
$$
\end{pro}

\begin{pro}
Prove that $3n^3-4n^2+13n=O(n^3)$
\begin{lenum}
\item
Using the definition of $O$.
\item
Using limits.
\end{lenum}
\end{pro}

\begin{pro}
Prove that $5n^2-7n=\Theta(n^2)$
\begin{lenum}
\item
Using the definition of $\Theta$ and/or Theorem~\ref{thm:theta}.
\item
Using limits.
\end{lenum}
\end{pro}

\begin{pro}
Prove that $n\log n=o(n^2)$.
\end{pro}

\begin{pro}
Prove that $\log(x^2+x)=\Theta(\log x)$.
\end{pro}

\begin{pro}
Prove that $\sqrt{5x^2+11x}=\Theta(x)$.
\end{pro}

\begin{pro}
Prove that $n^2=o(1.01^n)$.
\end{pro}

\begin{pro}
Consider the problem of computing the product of two matrices, $A$ and $B$, where 
$A$ is $l\times m$ and $B$ is $m\times n$.
\begin{lenum}
\item
Give an efficient algorithm to compute the product $A\times B$.
Assume you have a \verb+Matrix+ type with fields \verb+rows+ and \verb+columns+
that specify the number of rows/columns the matrix has.  Thus, you can call
\verb+A.rows+ to get the number of rows \verb+A+ has, for instance.
Also assume you can index a Matrix like an array.  Thus, \verb+A[i][j]+ accesses
the element in row $i$ and column $j$.
\item
Give the best and worst-case complexity of your algorithm.
\end{lenum}
\end{pro}


\begin{pro}
Give tight bounds for the {\bf best} and {\bf worst} case running times of each of the following algorithms in terms of the size of the input.
Assume $A.length=n$.  (Note: some of these are useless algorithms that do not do anything useful. Do not worry about what
they do, just how long it takes to do whatever it is they do.)
\begin{lenum}
\item
\begin{code}
void foo1(int n) {
    int foo = 0;
    for(int i = 0 ; i < n ; i++)
      foo += i;
}
\end{code}

\newpage

\item
\begin{code}
void blah(int n) {
    int blah = 0;
    for(int i = 0 ; i < sqrt(n) ; i++)
      blah += i;
}
\end{code}

\item
\begin{code}
void ferzle1(int a[], int n) {
    int ferzle = 0;
    for(int i = 0 ; i < n ; i++) {
        for(int j = 0 ; j < n ; j++) {
            ferzle += a[i]*a[j];
            if(ferzle==10000) {
                j=n;
            }
        }
    }
} 
\end{code}

\item
\begin{code}
void ferzle2(int n) {
    int ferzle = 0;
    for(int i = 0 ; i < n ; i++) {
        for(int j = i ; j < n ; j++) {
            ferzle += i*j;
        }
    }
}
\end{code}

\item
\begin{code}
void ferzle3(int a[], int n) {
    int ferzle = 0;
    for(int i = 0 ; i < n ; i++) {
        for(int j = 0 ; j < n ; j++) {
            ferzle += a[i]*a[j];
            if(ferzle==10000) {
                i=n;
            }
        }
    }
} 
\end{code}

\item
\begin{code}
void ferzle4(int a[], int n) {
    int ferzle = 0;
    for(int i = 0 ; i < n ; i++) {
        for(int j = 0 ; j < n ; j++) {
            ferzle += a[i]*a[j];
        }            
        if(ferzle==10000) {
            i=n;
        }
    }
} 
\end{code}

\newpage

\item
\begin{code}
void gruhop1(int n) {
    int gruhop = 0;
    for(int i = 0 ; i < n/2 ; i++) {
        for(int j = 0 ; j < n/2 ; j++) {
            gruhop += i*j;
        }
    }
}
\end{code}


\item
\begin{code}
int sumSomeStuff(int []A) {
    int sum=0;
    int i=0;
    while(i < A.length) {
        sum = sum + A[i];
        i++;
        if(sum > 100000) {
            i=A.length;
        }
    }
    return sum;
}
\end{code}

\item
\begin{code}
void gruhop2(int n) {
    int gruhop = 0;
    for(int i = 0 ; i < sqrt(n) ; i++) {
        for(int j = 0 ; j < n ; j++) {
            gruhop += i*j;
        }
    }
}
\end{code}

\item
\begin{code}
int doMoreStuff(int []A) {
    int sum=0;
    for(int i=0 ; i < A.length ; i++) {
        for(int j=0 ; j < A.length ; j++) {
            sum = sum + A[i]*A[j];
            if(sum==123) {
                j = A.length;
            }
        }
        for(int j=0 ; j < A.length ; j++) {
            sum = sum - A[j]*A[j];
        }
    }
    return sum;
}
\end{code}

\newpage

\item
\begin{code}
int sumTimesM(int []A) {
    int M = 100;
    int sum=0;
    for(int i=0 ; i < A.length ; i++) {
        for(int j=0 ; j < M ; j++) {
            sum = sum + A[j] + A[i];
            if(sum==123) {
                j = M;
            }
        }
    }
    return sum;
}
\end{code}

\item
\begin{code}
void foo2(int n,int m) {
    int foo = 0;
    for(int i = 0 ; i < n ; i++) 
        foo++; 
    for(int j = 0 ; j < m ; j++)
        foo++;
} 
\end{code}

\item
\begin{code}
void foo3(int n) { // Tricky one
    int foo = 0;
    for(int i = 1 ; sqrt(i)  <=  n ; i++)
        for(int j = 1 ; j <= i; j++) 
            doIt(j) // takes j steps;
}
\end{code}
%\begin{code}
%void foo2(int n) { // Tricky one
%    int foo = 0;
%    for(int i = 0 ; sqrt(i)  <  n ; i++)
%        for(int j = 0 ; j < i; j++) 
%            doIt(j) // takes j steps;
%}
%\end{code}

\item
\begin{code}
void HalfIt(int n) {
    while(n > 0) {
        n = n/2;
    }
}
\end{code}
\end{lenum}
\end{pro}

\begin{pro}
Using the code from Example~\ref{exa:retainAllOthers}, 
determine the complexity of the following method calls.
\begin{lenum}
\item
\verb+hs1.retainAll(al2)+, 
where \verb+hs1+ is a HashSet of size $n$ and \verb+al2+ is an ArrayList of size $m$.
\item
\verb+hs1.retainAll(ll2)+, 
where \verb+hs1+ is a HashSet of size $n$ and \verb+ll2+ is a LinkedList of size $m$.
\item
\verb+hs1.retainAll(ts2)+, 
where \verb+hs1+ is a HashSet of size $n$ and \verb+ts2+ is a TreeSet of size $m$.
\item
\verb+hs1.retainAll(hs2)+, 
where \verb+hs1+ is a HashSet of size $n$ and \verb+hs2+ is a HashSet of size $m$.


\item
\verb+ll1.retainAll(al2)+, 
where \verb+ll1+ is a LinkedList of size $n$ and \verb+al2+ is an ArrayList of size $m$.
\item
\verb+ll1.retainAll(ll2)+, 
where \verb+ll1+ is a LinkedList of size $n$ and \verb+ll2+ is a LinkedList of size $m$.
\item
\verb+ll1.retainAll(ts2)+, 
where \verb+ll1+ is a LinkedListof size $n$ and \verb+ts2+ is a TreeSet of size $m$.
\item
\verb+ll1.retainAll(hs2)+, 
where \verb+ll1+ is a LinkedList of size $n$ and \verb+hs2+ is a HashSet of size $m$.
\end{lenum}
\end{pro}

\newpage

\begin{pro}\index{sort!selection}\index{selection sort}
Consider the following two implementations of selection sort.
\begin{code}
void selectionSort(int a[],int n) {
    for (int i=0 ; i<n-1 ; i++) {
        int min = i;
        for (int j=i+1 ; j<n ; j++) {
            if(a[j] < a[min]) 
                min = j;
        }
        int temp = a[min];
        a[min] = a[i];
        a[i] = temp;
    }
}
\end{code}

\begin{code}
void selectionSort(List<Integer> a,int n) {
    for (int i=0 ; i<n-1 ; i++) {
        int min = i;
        for (int j=i+1 ; j<n ; j++) {
            if(a.get(j) < a.get(min)) 
                min = j;
        }
        int temp = a.get(min);
        a.set(min, a.get(i));
        a.set(i,temp);
    }
}
\end{code}
\begin{lenum}
\item
Give the worst-case complexity of the array version of selection sort (the first one).
\item
Give the worst-case complexity of the list version of selection sort (the second one) assuming the list is an {\em array-based} implementation (e.g. ArrayList).
\item
Give the worst-case complexity of the list version of selection sort (the second one) assuming the list is a {\em linked-list} implementation.
\item
Compare the three options.  Is one of them the clear choice to use?  Should any of them never be used? Explain.
\end{lenum}
\end{pro}


\begin{pro}
You need to choose data structures for two collections of data, $A$ and $B$, 
and the only thing you know is that the most common operation you will perform is
\verb+A.retainAll(B)+.
Given this, what are you best choices for $A$ and $B$?  Clearly justify your choices.
\end{pro}

\begin{pro}
In Java, a TreeMap is an implementation of the Map interface that uses a balanced binary search tree (a red-black tree) to store the keys and values.  In particular, the keys are used as keys in a BST with each key having an associated value.  TreeMaps have methods like\footnote{For simplicity we are ignoring generics here.  If you don't know what that means but you understand what this problem is saying, don't worry about it.} 
\verb+put(Object key, Object value)+ (add the key-value pair to the map), 
\verb+Object get(Object key)+ (returns the value
associated with the key), 
and \verb+ArrayList keySet()+ (returns an ArrayList of all the keys).  
As should be expected, \verb+put+ and \verb+get+ both take $\Theta(\log n)$ time.
You can assume that \verb+keySet+ takes $\Theta(n)$ time.

A method that might be useful on a TreeMap is \verb+ArrayList getAll(ArrayList keys)+ that
returns an ArrayList containing the values associated to the keys passed into the method. 
Consider the following implementation of this method.
\begin{code}
public ArrayList getAll(ArrayList keys) {
    ArrayList toReturn = new ArrayList();
    for (Object key : keySet()) {
        for (Object k : keys) {
            if (k.equals(key)) {
                toReturn.add(get(key));
            }
        }
    }
    return toReturn;
}
\end{code}
\begin{lenum}
\item
Does this method work properly?  Explain why it does or does not.
\item
What is the worst-case complexity of this method?
\item
Rewrite the method so that it is as efficient as possible and give the worst-case complexity of the new version.
\end{lenum}
\end{pro}

\begin{pro}
Consider the following implementation of binary search on a List.
\begin{code}
    int binarySearch(List A, int n, int val) {
       int left=0, right=n-1;
       while (right-left>=0) {
           int middle = (left+right)/2;
           if(val==A.get(middle)) 
               return middle;
           else if(val<A.get(middle)) 
               right=middle-1;
           else 
               left=middle+1; 
        }
       return -1;
    }
\end{code}
\begin{lenum}
\item
Give the worst-case complexity of this algorithm if $A$ is an array-based list (e.g., an ArrayList).
\item
Give the worst-case complexity of this algorithm if $A$ is linked list.
\item
Would it ever make sense to implement binary search on a linked list?  Explain.
\end{lenum}
\end{pro}

